\documentclass[12pt]{article}
\usepackage{comp4906}

\begin{document}
\thesisheader{Large Language Models as a Second Pass\texorpdfstring{\\}{ }%
  for Quantum Circuit Optimization}{Honours Thesis Literature Review}

\section{Circuit Optimization After Compilation}

A quantum circuit is a list of gates (operations on quantum bits, or qubits). Optimization finds a
less costly way to carry out the same operation. Two-qubit gates are a useful place to look for
savings because they are typically a major source of errors on quantum devices. For a given gate set
(the operations a circuit is allowed to use), two circuits can carry out the same operation with
different numbers of two-qubit gates. A compiler searches for the less expensive
implementation~\cite{ibm-opt}.

Qiskit (IBM's quantum software toolkit) can cancel redundant gates and use block resynthesis
(replacing groups of gates with equivalent implementations). Level 3 is its highest built-in
optimization setting. The result depends partly on which gates the compiler examines together. Two
identical controlled-NOT gates (CX gates, which flip one qubit when another is 1) cancel if the
compiler can bring them together. More involved reductions can require coordinating changes across a
larger part of the circuit. Minimizing CX count is NP-hard (a fast method guaranteed to find the
minimum would resolve a major open problem in computer science) even for circuits written in
Clifford+T (a standard fixed set of one- and two-qubit gates). Very small cases have useful exact
answers. For any operation on two qubits, with arbitrary one-qubit gates available, the minimum is
zero, one, two, or three CX gates and can be determined from the operation's
matrix~\cite{ibm-stages,wetering-amy,shende}.

Specialized optimizers can improve a circuit after a first compiler pass (one round of circuit
optimization). Quartz finds equivalent gate sequences and uses them as rewrite rules (ways of
replacing gates while preserving the operation). QUESO (a circuit optimizer that searches through such
rules) uses symbolic rules (patterns that cover many gate sequences) and beam search (keeping a
limited set of promising circuits at each step). The team behind quspect (another quantum compiler)
compared QUESO, GUOQ (an optimizer that combines rewrites with resynthesis), and Quartz on the same
circuits after Qiskit level 3 and routing (arranging qubit movements so gates can use a device's
connections). It reported further reductions in two-qubit gate counts, including about 7\% for QUESO
and 12\% for GUOQ on its evaluated circuit and device pairs. Two-qubit depth (the number of sequential
layers of two-qubit gates) increased in aggregate. That result establishes a useful comparison for a
large language model (LLM). An extra conventional search can already save gates on the compiled input,
and the accompanying change in depth can matter~\cite{quartz,queso,quspect}.

\section{What Language Models Have Shown}

Gu et al. gave Claude Mythos 5 a task specification and asked it to write a program that builds a
quantum circuit. The model could choose the algorithm and change how much temporary workspace it used.
An independent evaluator checked successive candidates and scored them by the product of their Toffoli
count (the number of gates that flip a target when both controls are 1) and peak qubit count (the most
qubits needed at once). The study shows that a language model can reorganize a quantum computation
when it can express and revise the algorithm directly. An existing compiled circuit provides a more
limited view of the computation. Its gate list can retain recognizable structure, such as repeated
carry steps (passing overflow between digits) in an adder (a circuit that adds numbers stored in
qubits). Jang et al. give a related reason to examine that structure. Their model used written design
guidance to produce circuits for specified functions, and the authors checked those circuits against
known input and output examples~\cite{gu,jang}.

StabilizerBench (a benchmark for quantum error-correction circuits) is the closest precedent for
rewriting a supplied circuit. In its optimization task, a model receives a circuit that prepares a
quantum state from zero and the stabilizers (operations that leave the required state unchanged) that
it must preserve. On 192 cases, Claude and GPT succeeded on about 73\% and 72\% of the circuits with
up to 15 verification calls and 900 seconds per case. Both succeeded substantially less often with the
same number of calls and a 300-second limit. A comparison with a conventional optimizer should hold
the time allowance fixed and record when a verified improvement becomes
available~\cite{stabilizerbench}.

The required behaviour also depends on the circuit's purpose. A circuit that prepares a state has one
prescribed input. A subroutine (a circuit called within a larger computation) may receive any input,
including a superposition (a combination of quantum states). A CX gate and an empty circuit agree on
the all-zero input. With the control qubit set to 1, only the CX gate flips the target. Checking the
full unitary (the matrix that describes the operation on every input) distinguishes the two. Comparing
these matrices up to global phase (one common complex multiplier of magnitude one) is the appropriate
check for replacing a standalone arithmetic or simulation block.

Saeedi et al. trained a model on circuits optimized by PyZX (a compiler that represents circuits as
diagrams and simplifies the diagrams). Among outputs that met their process-fidelity threshold (a
required degree of agreement between the old and new operations), the model matched PyZX's counts of T
gates (costly operations in error-corrected quantum computing). Correct outputs became much less
frequent as the gate lists grew. Petchartee's abstract reports reductions in gate count and depth for
QC-GPT (a model trained on quantum circuits from several software toolkits). The results motivate
tracking how often a model returns a valid circuit and how often that circuit improves on a
conventional result for the same input~\cite{drift,qcgpt}.

Crichigno and Flammia validate their checking procedure on solved cases and deliberately incorrect
candidates before evaluating their agents. That step is particularly useful in a gate-count
experiment. Deleting gates can lower the cost of a circuit while changing the operation it
performs~\cite{crichigno}.

\section{The Comparison This Study Makes}

Other benchmarks already compare multiple models on quantum compilation tasks. QuOptBench (a benchmark
for choosing a device and compiler settings) compares model choices with a search over the available
options. Qubrio (a compiler that schedules gates and qubit movements) compares estimates of execution
time and fidelity (how accurately the computation runs). Together with StabilizerBench, these studies
show that the contribution of a new benchmark comes from the decision it tests. Here that decision is
whether to give a circuit that has already passed through Qiskit to a model or to another conventional
optimizer. Both methods receive the same gate list, optimization objective, and time
allowance~\cite{quoptbench,qubrio}.

\researchquestion{When can LLMs improve an already optimized quantum circuit?}

There are two parts to the comparison. First, an extra conventional search establishes how many CX
gates the model would have to save to offer an advantage on the compiled input. Second, giving the
model both the original and optimized versions of the same operation shows whether its performance
changes after compilation. The two starting points should use the same gate notation. The paired
comparison should use the final CX count in each condition, since a longer starting circuit has more
gates available to remove. Correctness and elapsed time should be reported alongside gate savings,
with failed submissions counted even when the original circuit is retained.

The arithmetic cases are particularly informative if the model can recognize carry operations after
they have been expanded into smaller gates. Different adder designs help distinguish a recurring
result from one rewrite that happens to work on several widths of the same adder. Simulation and
variational circuits (templates with adjustable gates used in quantum algorithms) supply other forms
of repeated structure. An isolated improvement would justify looking closely at that construction.
Frequent ties with the conventional method would show how much of the apparent post-Qiskit saving is
already available from another optimizer. Comparing the outputs and their arrival times makes either
outcome interpretable.

\section{Extensions}

The most immediate improvement would be to compare any consistent model savings with a broader set of
compiler methods. tket (a general-purpose quantum compiler) can quickly simplify groups of gates. A
timed QUESO search gives a second reference. GUOQ or quspect would add further search strategies.
Cases with known CX minima are also useful controls. If an exact minimum is available, a tie can be
interpreted against that minimum. On the main workload, a tie is evidence about the methods and time
tested~\cite{tket,quspect,shende,shaik-vandepol}.

\begin{itemize}
  \item \textbf{Ask for edits to the input circuit.} Have the model identify gates to replace and let
  the test program copy the rest. Comparing edits with complete gate lists on the same cases would
  help determine whether failures arise from having to reproduce long stretches of unchanged
  gates~\cite{drift}.
  \item \textbf{Test whether knowing the circuit's purpose helps.} Give the model a short description,
  such as the kind of addition the circuit performs. Compare this with a conventional synthesis method
  that can use the same description. This separates the value of identifying the operation from the
  value of optimizing its gate list~\cite{gu}.
  \item \textbf{Ask the model to write a reusable compiler pass.} The model would supply a rule or a
  search procedure once, and the study would test it on independent circuit families. QuantumEvo (a
  system that designs algorithms for building quantum circuits) uses this approach to design a
  reusable heuristic (a practical decision rule) for building circuits whose operations can be undone.
  Reuse matters if an effective per-circuit rewrite costs too much to repeat at compilation
  time~\cite{quantumevo}.
  \item \textbf{Measure savings after device mapping.} Apply the same routing and native gate set (the
  gates available on the device) to the original and rewritten circuits. This would establish whether
  the logical gate savings survive hardware constraints and provide a basis for measuring execution
  time or success on a device~\cite{qubrio}.
\end{itemize}


\thesisreferences
\end{document}
